\section{The One Substance}
\label{sec:thesis}

Start with the boldest claim, because everything else is its consequence. Reality is one
substance, and that substance is experiential. The smallest unit is the vibe, a tiny ripple
of feeling. It is not a particle with an experience painted on it. It is experience, the
thing itself, however faint, at the smallest scale.

\begin{principle}[Monism]
There is one kind of entity, the vibe. Every other thing, space, time, matter, force, mind,
is a pattern of vibes and their relations, not a separate ingredient.
\end{principle}

This turns the usual picture upside down. The standard story begins with dead matter and
hopes that, arranged cleverly enough, it will somehow produce consciousness, and no one can
say how. Vibe Theory begins with experience as the ground and treats dead-seeming matter as
what experience looks like once it has settled into stable, shared, repeatable patterns.
Matter, on this view, is the late arrival, not the foundation. On this view, too, mind is not
produced by the brain. The brain is taken to be a particularly intricate, particularly
integrated knot of the one experiential field.

It helps to see the substance as ranging between two conditions.

\begin{itemize}
\item The \emph{lawful} is whatever has settled into regular, stable, structured pattern,
holding its form steadily.
\item The \emph{wild} is whatever stays fluid, free, and unsettled.
\end{itemize}

This is a matter of how settled a pattern is, and it cuts across every domain rather than
naming a kind of thing. Lawful does not mean physical. The physical world is the most familiar
lawful region, the part of the field that has crystallized into a shared, stable pattern the
same for everyone, with stone and light and the laws of motion among the dreams the universe
has dreamt so consistently that they have become solid. But a disciplined mind or a settled
inner structure can be highly lawful in its own domain, and a spiritual order, if there is one,
could be lawful too, regular and calm and structured without being physical at all. Equally,
every domain has its wild side, the part still fluid and unmade, with much of imagination,
dreaming, and raw feeling among it. So the lawful and the wild are conditions a pattern can be
in, found across the physical, the mental, and whatever else, not another name for outer versus
inner. They are one ocean, part frozen into reliable ice and part still moving water, and ice
and water can form anywhere in it. There is no second substance dividing them, only one
substance in different conditions.

The discipline this paper imposes on itself, and the reason the picture is testable rather
than merely poetic, is that the physical world, the shared lawful pattern, must be
\emph{derived}, not assumed. It is not enough to say that physics could come out of experience.
The program has been to build the substrate, run it, and check that space, matter, gravity, and
the quantum actually emerge, with numbers that could have come out wrong. Most of this paper is
the report of that emergence in plain terms. The felt interior is where the model can build and
measure the structure but cannot, from the outside, settle whether the structure is the same
as the feeling. That boundary, the hard problem at the center of any experience-first view,
is named honestly when we reach it.
